% Options for packages loaded elsewhere
\PassOptionsToPackage{unicode}{hyperref}
\PassOptionsToPackage{hyphens}{url}
\documentclass[
]{article}
\usepackage{xcolor}
\usepackage[margin=1in]{geometry}
\usepackage{amsmath,amssymb}
\setcounter{secnumdepth}{-\maxdimen} % remove section numbering
\usepackage{iftex}
\ifPDFTeX
  \usepackage[T1]{fontenc}
  \usepackage[utf8]{inputenc}
  \usepackage{textcomp} % provide euro and other symbols
\else % if luatex or xetex
  \usepackage{unicode-math} % this also loads fontspec
  \defaultfontfeatures{Scale=MatchLowercase}
  \defaultfontfeatures[\rmfamily]{Ligatures=TeX,Scale=1}
\fi
\usepackage{lmodern}
\ifPDFTeX\else
  % xetex/luatex font selection
\fi
% Use upquote if available, for straight quotes in verbatim environments
\IfFileExists{upquote.sty}{\usepackage{upquote}}{}
\IfFileExists{microtype.sty}{% use microtype if available
  \usepackage[]{microtype}
  \UseMicrotypeSet[protrusion]{basicmath} % disable protrusion for tt fonts
}{}
\makeatletter
\@ifundefined{KOMAClassName}{% if non-KOMA class
  \IfFileExists{parskip.sty}{%
    \usepackage{parskip}
  }{% else
    \setlength{\parindent}{0pt}
    \setlength{\parskip}{6pt plus 2pt minus 1pt}}
}{% if KOMA class
  \KOMAoptions{parskip=half}}
\makeatother
\usepackage{longtable,booktabs,array}
\usepackage{calc} % for calculating minipage widths
% Correct order of tables after \paragraph or \subparagraph
\usepackage{etoolbox}
\makeatletter
\patchcmd\longtable{\par}{\if@noskipsec\mbox{}\fi\par}{}{}
\makeatother
% Allow footnotes in longtable head/foot
\IfFileExists{footnotehyper.sty}{\usepackage{footnotehyper}}{\usepackage{footnote}}
\makesavenoteenv{longtable}
\usepackage{graphicx}
\makeatletter
\newsavebox\pandoc@box
\newcommand*\pandocbounded[1]{% scales image to fit in text height/width
  \sbox\pandoc@box{#1}%
  \Gscale@div\@tempa{\textheight}{\dimexpr\ht\pandoc@box+\dp\pandoc@box\relax}%
  \Gscale@div\@tempb{\linewidth}{\wd\pandoc@box}%
  \ifdim\@tempb\p@<\@tempa\p@\let\@tempa\@tempb\fi% select the smaller of both
  \ifdim\@tempa\p@<\p@\scalebox{\@tempa}{\usebox\pandoc@box}%
  \else\usebox{\pandoc@box}%
  \fi%
}
% Set default figure placement to htbp
\def\fps@figure{htbp}
\makeatother
\setlength{\emergencystretch}{3em} % prevent overfull lines
\providecommand{\tightlist}{%
  \setlength{\itemsep}{0pt}\setlength{\parskip}{0pt}}
\usepackage{graphicx}
\usepackage{float}
\usepackage{amsmath}
\usepackage{bookmark}
\IfFileExists{xurl.sty}{\usepackage{xurl}}{} % add URL line breaks if available
\urlstyle{same}
\hypersetup{
  hidelinks,
  pdfcreator={LaTeX via pandoc}}

\author{}
\date{\vspace{-2.5em}}

\begin{document}

\section{Question}\label{question}

Gustavo, coordinador académico, registró los puntajes del simulacro
ICFES. La tabla muestra los datos ordenados de menor a mayor.

\begin{longtable}[]{@{}ccc@{}}
\toprule\noalign{}
Estudiante & puntajes (puntos) & Cuartil \\
\midrule\noalign{}
\endhead
\bottomrule\noalign{}
\endlastfoot
1 & 257 & Mín \\
2 & 263 & \\
3 & 267 & Q1 \\
4 & 289 & \\
5 & 333 & \\
6 & 334 & Q2 \\
7 & 338 & \\
8 & 347 & \\
9 & 374 & Q3 \\
10 & 399 & \\
11 & 400 & Máx \\
\end{longtable}

Un estudiante de Gustavo intentó construir el diagrama de caja
correspondiente a estos datos.

\textbf{¿Cuál de los siguientes diagramas de caja representa
CORRECTAMENTE la información de la tabla?}

\subsection{Answerlist}\label{answerlist}

\begin{itemize}
\tightlist
\item
  \includegraphics[width=0.6\linewidth,height=\textheight,keepaspectratio]{diagrama_a.png}
\item
  \includegraphics[width=0.6\linewidth,height=\textheight,keepaspectratio]{diagrama_b.png}
\item
  \includegraphics[width=0.6\linewidth,height=\textheight,keepaspectratio]{diagrama_c.png}
\item
  \includegraphics[width=0.6\linewidth,height=\textheight,keepaspectratio]{diagrama_d.png}
\end{itemize}

\section{Solution}\label{solution}

\subsubsection{Información de la
Pregunta}\label{informaciuxf3n-de-la-pregunta}

\begin{longtable}[]{@{}
  >{\raggedright\arraybackslash}p{(\linewidth - 2\tabcolsep) * \real{0.4091}}
  >{\raggedright\arraybackslash}p{(\linewidth - 2\tabcolsep) * \real{0.5909}}@{}}
\toprule\noalign{}
\begin{minipage}[b]{\linewidth}\raggedright
Aspecto
\end{minipage} & \begin{minipage}[b]{\linewidth}\raggedright
Descripción
\end{minipage} \\
\midrule\noalign{}
\endhead
\bottomrule\noalign{}
\endlastfoot
\textbf{Competencia} & Interpretación y representación \\
\textbf{Componente} & Aleatorio \\
\textbf{Afirmación} & Interpreta representaciones gráficas de datos
estadísticos \\
\textbf{Evidencia} & Identifica correspondencia entre tabla de datos y
diagrama de caja \\
\textbf{Tarea} & Identificar el diagrama de caja que representa
correctamente los cuartiles de un conjunto de datos ordenados \\
\textbf{Nivel de dificultad} & N2 \\
\end{longtable}

\subsubsection{¿Qué evalúa esta
pregunta?}\label{quuxe9-evaluxfaa-esta-pregunta}

Esta pregunta evalúa la capacidad del estudiante para interpretar tablas
de datos ordenados y relacionarlas con su representación gráfica
correcta mediante un diagrama de caja y bigotes. El estudiante debe
identificar los cinco valores característicos (mínimo, Q1, mediana, Q3,
máximo) y verificar cuál diagrama los representa correctamente.

\subsubsection{Respuesta Correcta: Opción
D}\label{respuesta-correcta-opciuxf3n-d}

\textbf{Justificación matemática:}

Para construir correctamente un diagrama de caja con \(n = 11\) datos
ordenados, debemos calcular las posiciones de los cuartiles:

\textbf{Paso 1:} Calcular la posición de cada cuartil
\[Q_1 = \text{dato en posición } \frac{n+1}{4} = \frac{11+1}{4} = 3\]
\[Q_2 = \text{dato en posición } \frac{n+1}{2} = \frac{11+1}{2} = 6\]
\[Q_3 = \text{dato en posición } \frac{3(n+1)}{4} = \frac{3 \times 12}{4} = 9\]

\textbf{Paso 2:} Leer los valores correspondientes de la tabla ordenada

\begin{longtable}[]{@{}lll@{}}
\toprule\noalign{}
Medida & Posición & Valor \\
\midrule\noalign{}
\endhead
\bottomrule\noalign{}
\endlastfoot
Mínimo & 1 & 257 puntos \\
Q1 & 3 & 267 puntos \\
Q2 (Mediana) & 6 & 334 puntos \\
Q3 & 9 & 374 puntos \\
Máximo & 11 & 400 puntos \\
\end{longtable}

\textbf{Por lo tanto}, el diagrama correcto es el que muestra
exactamente estos cinco valores: Mín = 257, Q1 = 267, Q2 = 334, Q3 =
374, Máx = 400.

\includegraphics[width=0.7\linewidth,height=\textheight,keepaspectratio]{diagrama_d.png}

\subsubsection{Opciones No Válidas}\label{opciones-no-vuxe1lidas}

Las opciones incorrectas contienen errores conceptuales comunes en el
cálculo de cuartiles:

\textbf{Error del diagrama con EST-BOX-03:}

Es posible que los estudiantes que eligen esta opción ubicó el primer
cuartil (q1) donde debería estar el tercer cuartil (q3) y viceversa.
este error indica confusión sobre qué porcentaje de datos queda debajo
de cada cuartil. Este error se presenta cuando confunde el significado
de q1 (25\%) y q3 (75\%) en la distribución. Para evitar este error, el
estudiante debe usar las fórmulas correctas de posición: Q1 en posición
(n+1)/4, Q2 en (n+1)/2, Q3 en 3(n+1)/4.

\textbf{Otros errores posibles en los diagramas incorrectos:}

\begin{itemize}
\item
  \textbf{Cuartiles intercambiados (EST-BOX-03):} Es posible que los
  estudiantes confundan Q1 con Q3, ubicando el tercer cuartil donde
  debería ir el primero. Para evitar este error, verificar siempre que
  Q1 \textless{} Q2 \textless{} Q3.
\item
  \textbf{Q1 mal calculado (EST-BOX-02):} Es posible que los estudiantes
  ubiquen Q1 muy cerca del mínimo. Para evitar este error, recordar que
  Q1 está en la posición 3 para n=11 datos.
\item
  \textbf{Mediana como promedio (EST-BOX-04):} Es posible que los
  estudiantes confundan la mediana con el promedio aritmético. Para
  evitar este error, recordar que la mediana es el valor central, no un
  cálculo aritmético.
\end{itemize}

\subsubsection{Reflexión
Metacognitiva}\label{reflexiuxf3n-metacognitiva}

La caja del diagrama contiene el 50\% central de los datos, entre Q1 y
Q3.

\textbf{Estrategias para evitar errores comunes:}

\begin{enumerate}
\def\labelenumi{\arabic{enumi}.}
\tightlist
\item
  \textbf{Ordenar primero}: Los cuartiles solo se calculan sobre datos
  ordenados de menor a mayor
\item
  \textbf{Usar fórmulas de posición}: Para n datos, Q1 está en posición
  (n+1)/4, Q2 en (n+1)/2, Q3 en 3(n+1)/4
\item
  \textbf{Verificar coherencia}: Siempre debe cumplirse
  \(\text{Mín} \leq Q_1 \leq Q_2 \leq Q_3 \leq \text{Máx}\)
\end{enumerate}

\subsection{Answerlist}\label{answerlist-1}

\begin{itemize}
\tightlist
\item
  Incorrecto. Este diagrama tiene errores en la ubicación de los
  cuartiles.
\item
  Incorrecto. Este diagrama tiene errores en la ubicación de los
  cuartiles.
\item
  Incorrecto. Este diagrama tiene errores en la ubicación de los
  cuartiles.
\item
  CORRECTO. Este diagrama representa correctamente los cuartiles de los
  datos.
\end{itemize}

\section{Meta-information}\label{meta-information}

exname:
diagrama\_caja\_estaturas\_metacognitivo\_interpretacion\_n2\_schoice\_v1
extype: schoice exsolution: 0001 exshuffle: FALSE extol: 0.01 exsection:
Estadistica\textbar Medidas-Posicion\textbar Diagramas-Caja

exextra{[}Type{]}: SCHOICE exextra{[}Program{]}: R
exextra{[}Language{]}: es

exextra{[}Competencia{]}: Interpretacion exextra{[}Componente{]}:
Aleatorio exextra{[}Pensamiento{]}: Aleatorio exextra{[}Afirmacion{]}:
Interpreta representaciones graficas de datos estadisticos
exextra{[}Evidencia{]}: Identifica correspondencia entre tabla de datos
y diagrama de caja exextra{[}Nivel{]}: 2 exextra{[}Contexto{]}: Escolar

exextra{[}DOK{]}: 2 exextra{[}Bloom{]}: Analizar exextra{[}SOLO{]}:
Relacional

exextra{[}TipoMetacognicion{]}: identificacion\_grafico\_correcto

\end{document}
